% ============================================================================
% Remerciements
% ============================================================================
\chapter*{Remerciements}
\addcontentsline{toc}{chapter}{Remerciements}
\markboth{Remerciements}{Remerciements}

\vspace{0.5cm}

Nous tenons à exprimer notre profonde gratitude envers toutes les personnes
qui ont contribué, de près ou de loin, à la réalisation de ce projet.

Nos remerciements les plus sincères vont en premier lieu à notre encadrant,
\textbf{Monsieur Ezzati}, pour bien plus qu'un simple encadrement. Sa rigueur,
ses encouragements constants et sa capacité à nous pousser à dépasser nos
limites ont été le moteur de tout ce travail. Il nous a accompagnés avec
bienveillance tout au long du semestre, toujours disponible pour orienter
nos réflexions et corriger notre cap. Nous sommes véritablement fiers de ce
que nous avons accompli, et cela n'aurait pas été possible sans son soutien.
Qu'il reçoive ici l'expression de toute notre reconnaissance et de notre
admiration.

Un merci sincère à nos familles respectives pour leur soutien moral
indéfectible, leur patience et leur compréhension pendant les longues
nuits de développement et de rédaction.

Nous remercions également la communauté \textit{open source} mondiale,
dont les contributions ont rendu possibles des technologies telles que
PostgreSQL, TimescaleDB, NATS, FastAPI, Ollama, React et Docker — autant
de briques fondamentales de la plateforme \caimp{}. Ces projets, librement
partagés, incarnent l'idéal de la coopération intellectuelle.

Enfin, nous remercions nos camarades de promotion pour les échanges
fructueux, les relectures et l'esprit de camaraderie qui ont animé cette
année universitaire.

\bigskip
\begin{flushright}
\textit{CHAHIDY Rim, SAMDI Wiam, TAÏK Youssef}\\
\textit{Settat, 2026}
\end{flushright}
